\documentclass[10pt]{article}

\usepackage[utf8]{inputenc}

\usepackage[T1]{fontenc}

\usepackage{amsmath}

\usepackage{amsfonts}

\usepackage{amssymb}

\usepackage[version=4]{mhchem}

\usepackage{stmaryrd}



\begin{document}

ются евклидовыми или унитарными пространствами и получение оценок типа (3) возможно на основе относительно простых средств. Напр., пусть рассматривается линейная двухслойная разностная схема вида



\[

\left.\begin{array}{c}

A_{1} u^{n+1}=A_{0} u^{n}+\tau f^{n+1}, n \geqslant 0, \\

u^{0}=\varphi

\end{array}\right\}

\]



где векторы $\varphi$ и $f^{n+1}$ определяются начальным условием и правой тастью уравнения, а операторы $A_{1}$ и $A_{0}$ в евклидовом пространстве $H$ таковы, тто $\left\|A_{1}^{-1}\right\| \leqslant$ $\leqslant C_{0},\left\|A_{1}^{-1} A_{0}\right\| \leqslant 1+C_{1} \tau$, где неотрицательные константы $C_{0}$ и $C_{1}$ не зависят от сетки. Тогда для решения (6) справедлива априорная оценка



\[

\max _{n, t \in \omega_{t}}\left\|u^{n}\right\|_{H} \leqslant e^{C_{1} T}\left\{\|\varphi\|_{H}+C_{0} \sum_{n, t_{n+1} \in \omega_{t}} \tau\left\|f^{n+1}\right\| H\right\} \text {, }

\]



Весьма часто анализ таких схем проводится после записи их в канонич. виде



\[

B\left(\frac{u^{n+1}-u^{n}}{\tau}\right)+A u^{n}=f^{n+1}

\]



на основе изучөния свойств оператора перехода $R \equiv$ $\equiv E-\tau B^{-1} A \quad(E-$ тождественный оператор) в предположении, что имеется нек-рая относительно простая информация типа операторных неравенств об операторах $B$ и $A$ в евклидовом пространстве $H$. Напр., если $B=B^{*}>0, A=A^{*}$ и $B \geqslant \tau(1+\rho)^{-1} A$, то (см. [3], [6]) существует константа $C \equiv C(T, \rho)$ такая, что



\[

\left\|u^{n+1}\right\|_{B} \leqslant C\left(\left\|u^{0}\right\|_{B}+\tau \sum_{k=0}^{n}\left\|f^{k+1}\right\|_{B-1}\right),

\]



где $\left\|u^{n}\right\|_{B} \equiv\left(B u^{n}, u^{n}\right)_{H}^{1 / 2}$. Подобные результаты полутены для достаточно широкого круга разностных схем, включая трехслойные п нек-рые многослойные схемы (см. [6]). При этом изучены ІІ нек-рые частные случаи устойчивости (устойчивость по начальным данным, устойчивость по правой части) и их взаимоотношения.



Имеются нек-рые результаты, связанные с изучением необходимых условий подобной устойчивости или близких к ним (см. [3], [6]). Использование энергетич. неравенств (см. [4], [5]) вместо (8) позволяет при родственных условиях получить оценку типа



\[

\begin{aligned}

& \left\|u^{n+1}\right\|_{B}^{2}+\tau \sum_{k=0}^{n}\left\|u^{k+1}\right\|_{H}^{2} \leqslant \\

\leqslant & C_{1}\left(\left\|u^{0}\right\|_{B}^{2}+\tau \sum_{k=0}^{n}\left\|f^{k+1}\right\|_{B-1}^{2}\right),

\end{aligned}

\]



приводящую к устойчивости в несколько более сильной норме для $u_{h}$ и переходящей в пределе в оценку, часто встречающуюся в теории әволюционных уравнений. Подобные оценки также получены для весьма широкого круга схем (см. [4], [5], [13]).



При изучении У. р. с. выделяют условно устойчивые разностные схемы типа явных схем для уравнения теплопроводности, в к-рых устойчивость имеется лишь при ограничениях типа $\tau\left(h_{s}\right)^{-2} \leqslant x_{0}$, и схемы абсолютно устойчивые, в к-рых шаги по времени и іго пространственным переменным могут меняться независимо друг от друга, не нарушая устойчивости. Схемы последнего типа часто являются предпочтительными, если они не требую' решения сложных систем на каждом шаге. If таким әкономичным разностным схемам для многомерных задач относятся неявные схемы переменных направлений, схемы расщепления, схемы с расщепляющимся оператором и аддитивные схемы (см. [3] - [6]).



Теоремы устойчивости и оценки типа (3), (9) находят применение и в случае, когда погрешность аппроксимации и оценка (5) не рассматриваются, а строятся соответствующие восполнения решений сеточных задач и устанавливается на основе теорем компактности сходимость к решению исходной задачи (см. [4], [5]). Использование различных априорных оценок и упомянутого принципа компактности особенно характерно для сложных нелинейных задач, в к-рых решение может быть и неединственно, а сходимость устанавливается лишь к нек-рому решению исходной задачи. Иногда изучение нелинейных задач математич. физики по причине их сложности вообще заменяется изучением их линеаризацией, а для разностных схем обращается особое внимание на справедливость сеточных аналогов важнейших физич. законов сохранения (см. [8]). Для слабо же нелинейных задач изучение корректности разностных схем часто проводится с достаточной полнотой, характерной для линейного случая (см. [5] - [7] и Нелинейная краевая задача; числөнные методы решения).



В случае задач Коши для систем обыкновенных дифференциальных уравнений изучение устойчивости разностных схем часто сводится в модельных ситуациях к изучению корней характеристич. уравнения (cM. [2], [14], [15]).



Jum.: [1] Р я б е н в к и й В. С., Ф и ли и пов в А. Ф., Об устойчивости разностных уравнений, М., 1956, [2] Б е р ез и н И. С., अ З и д к о в Н. П., Методы вычислений, т. 2, 2 изд., ностные схемы, 2 изд., М., 1977, [4] ЈI а д ы ж е н с К а я О. А., Краевые задачи математической физики, М., 1973; [5] Д ь я к он о в Е. Г., Разностные методы решения краевых задач, в. $1-2$,

М. $1971-72 ;[6]$ С а м а р с и й А. А., Г у л и н А. В., Устойчивость разностных схем, М., 1973; [7]'С а м а р с к и й А. А., тойдасть разностных схем, М., уравнений, М., 1976; [ 8] С а м а р с к и й А. А., П о п о в Ю. П., Разностные методы решения задач газовой динамики, 2 изд., росы динамики вязкой несжимаемой жидкости, 2 изд., М., 1970; [10] О г а н е с я Л. А., Р у х о в е І Л. А., Вариаџионноразностные методы решения әллиптических уравнений, Ер., 1979; [11] М и х л и н С. Г. Численная реализация вариационных методов, М., 1966; [12] С т р е н г Г., Ф и к с Д Ж., Тео-

рия метода конечных элементов, пер. с англ., М., $1977 ;$ [13] рия метода конечных әлементов, пер. с англ., М., 1977; [13] с. 27-35; [14] Ба а х в а лі о в Н. С., Численные методы, 2 изд., ч., ких систем, М., 1979.



УСТОЙЧИВОСТЬ УПРУ ИИХ СИСТЕМ - 1) У. У. с.свойство упругих систем (упругих тел или совокупностей взаимодействующих упругих тел) мало отклоняться от состояния равновесия (движения) при достаточно малых возмущающих воздействиях. Роль возмущающих воздействий играют флуктуации внешних сил, отклонения от идеальной геометрич. формы, дефекты материала и т. п.



\begin{enumerate}

  \setcounter{enumi}{1}

  \item У. у. с. - раздел м е х а н и и и д е фо р м ир у е м го т в е рд о г о те л а, предметом к-рого является изучение У. У. с. В более широкой трактовке этот разделг включает изучение устойчивости вязкоупругих, упруго-пластических и др. деформируемых систем; часто употребляют также термин у с т о йч и в ость д е ф о р и и у е мых систем.

\end{enumerate}



Понятие У. У. с. тесно связано с общим понятием устойчивости движения, в частности с понятием устойчивости по Ляпунову. Центральная задача теории У. у. с.- нахождение области в пространстве параметров системы и внешних воздействий, в пределах к-рой рассматриваемое состояние равновесия (движения) остается устойчивым. Поверхность, ограничивающая область устойчивости, наз. к р и т и ч е с к о й п ов е р х о с т ю. Часто воздействие на упругую систему задают с точностью до одного параметра $\beta$. Без ограничения общности можно принять, что $0 \leqslant \beta<\infty$, причем при $\beta=0$ имеет место устойчивость. Нижняя грань $\beta_{*}$ значений параметра $\beta$, при к-рых исследуемое равновесие (движение) перестает быть устойчивым, наз. к р и т и ч е с к и м п а р а м е т р ол. Задачи





\end{document}